% Ilustraciones conceptuales de Lumina Retail.
% Requiere graphicx; cargar este archivo una vez en el preámbulo.
% Las decisiones pertenecen a la Arquitectura de Solución y al diseño aprobado.
% Las escenas ilustran conceptos de negocio; no especifican contratos o servicios.

\providecommand{\LuminaAnchoFigura}{0.55\linewidth}
\providecommand{\LuminaRutaFiguras}{../report-resources/figures/conceptuales/}

\newcommand{\LuminaFigura}[4]{%
  \begin{figure}[htbp]
    \centering
    \includegraphics[width=\LuminaAnchoFigura,height=4.8cm,keepaspectratio]{\LuminaRutaFiguras#1}%
    \caption[#2]{#3}%
    \label{#4}%
  \end{figure}%
}

% Portada: usar dentro de titlepage, fuera de figure y sin caption.
\newcommand{\LuminaPortada}{%
  \par\begingroup\centering
  \includegraphics[width=11cm,height=7cm,keepaspectratio]{\LuminaRutaFiguras00_portada_arquitectura_solucion.png}\par
  \endgroup
}

\newcommand{\LuminaSilos}{%
  \LuminaFigura{01_pedido_silos.png}{El pedido que da vueltas}%
  {Inventarios e información aislados dificultan atender un pedido aun cuando el producto existe en la red.}%
  {fig:concepto-silos}%
}

\newcommand{\LuminaFuentes}{%
  \LuminaFigura{02_fuentes_hibridas.png}{Dos entradas, una operación}%
  {Históricos de referencia y eventos sintéticos aportan funciones complementarias; la integración conserva el origen y utiliza correspondencias validadas.}%
  {fig:concepto-fuentes}%
}

\newcommand{\LuminaCapas}{%
  \begin{figure}[htbp]
    \centering
    \begin{minipage}{\LuminaAnchoFigura}
      \centering
      \includegraphics[width=\linewidth,height=4.8cm,keepaspectratio]{\LuminaRutaFiguras03_bronze_silver_gold_tres_mesas.png}\par
      \smallskip
      \parbox[t]{0.32\linewidth}{\centering\small\textbf{Bronze}\\Conservar}%
      \hfill
      \parbox[t]{0.32\linewidth}{\centering\small\textbf{Silver}\\Preparar}%
      \hfill
      \parbox[t]{0.32\linewidth}{\centering\small\textbf{Gold}\\Servir información}%
    \end{minipage}
    \caption[Tres mesas de trabajo]{Tres responsabilidades de la plataforma: preservar los originales, preparar datos consistentes y producir información útil para el negocio.}%
    \label{fig:concepto-capas}%
  \end{figure}%
}

\newcommand{\LuminaContinuidad}{%
  \LuminaFigura{04_inventario_indisponible.png}{La tienda perdió la señal}%
  {Ante una fuente de inventario indisponible, la vigencia del dato y la verificación operativa condicionan la confirmación del pedido.}%
  {fig:concepto-continuidad}%
}

\newcommand{\LuminaGobierno}{%
  \LuminaFigura{05_pasaporte_dato.png}{El pasaporte del dato}%
  {Cada producto de información debe conservar trazabilidad de su origen, sus transformaciones y sus controles de acceso.}%
  {fig:concepto-gobierno}%
}

\newcommand{\LuminaValor}{%
  \LuminaFigura{06_alerta_valor.png}{De la alerta al valor}%
  {El valor esperado se verifica relacionando una señal con la actuación realizada y su resultado operativo.}%
  {fig:concepto-valor}%
}
